% Modelo de livro da Editora CSVET (LaTeX). Compile com:  latexmk livro.tex
% Opções: [16x23] (padrão) ou [21x28]; [coletanea] para obras organizadas; [comercial] para obra sem licença aberta.
\documentclass[16x23]{csvet-livro}

\title{Título do Livro}
\subtitulo{Subtítulo do livro}
\autores{Nome do Autor\enspace·\enspace Nome da Autora}
\orientacao{Prof. Dr. Nome do Orientador}          % apague em obras que não derivam de trabalho acadêmico
\edicao{1ª edição}
\ano{2026}
\isbn{ISBN 978-65-00-00000-0 (PDF)\quad·\quad ISBN 978-65-00-00000-0 (impresso)}
\doi{}                                              % ex.: https://doi.org/10.00000/xxxx
\notadeorigem{Esta obra é uma versão revista e adaptada do Trabalho de Conclusão de Curso (Graduação em Medicina Veterinária)
  de Nome do Autor, aprovado em 10/12/2025 na Universidade Exemplo, sob orientação de Prof. Dr. Nome do Orientador.
  Banca examinadora: Prof.ª Dr.ª Nome (Instituição) e Prof. Me. Nome (Instituição). Orientação e banca não são autoras desta edição.}
% Obras de membros da editora (veja a política de conflito de interesses):
% \notaeditorial{Os autores integram a Editora CSVET. A avaliação desta obra foi conduzida de forma independente por Prof.ª Dr.ª Nome
%   (Conselho Editorial), sem participação dos autores na escolha dos pareceristas nem na decisão editorial.}

\begin{document}

\paginasiniciais{%
  \begin{expediente}
    \papel{Editor-chefe}{Prof. Dr. Atilio Sersun Calefi – Universidade Cruzeiro do Sul, Brasil}
    \papel{Conselho Editorial}{Prof.ª Dr.ª Nome Sobrenome – Instituição, País\newline Prof. Dr. Nome Sobrenome – Instituição, País\newline
      Prof.ª Dr.ª Nome Sobrenome – Instituição, País\newline Prof. Dr. Nome Sobrenome – Instituição, País}
    \papel{Pareceristas desta obra}{Prof. Dr. Nome Sobrenome – Instituição, País\newline Prof.ª Dr.ª Nome Sobrenome – Instituição, País\newline
      {\itshape (avaliação duplo-anônima; nomes publicados após a aprovação, com autorização)}}
    \papel{Revisão de texto}{Me. Nome Sobrenome – Letras, Instituição}
    \papel{Normalização}{Nome Sobrenome – Bibliotecário(a), CRB-0/0000}
    \papel{Diagramação}{Nome Sobrenome}
    \papel{Capa}{Nome Sobrenome (imagem: fonte e licença)}
  \end{expediente}}

\chapter*{Apresentação}
\addcontentsline{toc}{chapter}{Apresentação}
Texto da apresentação (organizadores) ou do prefácio (convidado): objetivo do livro, público, organização dos capítulos e como usar a obra.

\mainmatter
% Em coletâneas, use \capitulo[Autores do capítulo]{Título}; em obras autorais, \capitulo{Título}.
\capitulo[Nome do Autor · Nome da Autora]{Título do primeiro capítulo}

\begin{resumo}
Até 150 palavras (coletâneas): objetivo, conteúdo principal e contribuição do capítulo.
\end{resumo}
\palavraschave{termo; termo; termo.}

\section{Seção}
Texto corrido em Source Serif 4, corpo 10,5 pt sobre 14,5 pt, justificado, com recuo de 5 mm na primeira linha. A
hemogasometria\index{Hemogasometria} é usada na rotina de cães\index{Cães} e gatos\index{Gatos}. Para incluir um termo no
índice remissivo, escreva o comando \texttt{\textbackslash index\{Termo\}} logo depois dele; para subentradas, use
\texttt{\textbackslash index\{Acidose!metabólica\}}.

Esta mancha gráfica e esta entrelinha dão de 2.000 a 2.200 caracteres por página no formato 16×23 cm, com boa leitura na tela
e no papel. Títulos em Roboto, na cor da marca. Capítulos sempre abrem em página ímpar, sem cabeçalho.

\begin{citacao}
Citação direta com mais de três linhas, recuada 12 mm, em corpo menor e sem aspas, seguida da indicação da fonte com a página
(Autor, 2024, p. 00).
\end{citacao}

\begin{figure}[htb]
  \centering
  \includegraphics[width=\textwidth]{figura-exemplo}
  \caption{Legenda descritiva da figura, que permita entendê-la sem ler o texto.}
  \fonte{os autores (2026).}
\end{figure}

\begin{table}[htb]
  \caption{Título da tabela, acima dela.}
  \centering\small
  \begin{tabular}{lcc}
    \toprule
    Variável & Grupo A & Grupo B \\
    \midrule
    pH & 7,38 ± 0,02 & 7,21 ± 0,05 \\
    HCO$_3^-$ (mmol/L) & 22 ± 2 & 14 ± 3 \\
    \bottomrule
  \end{tabular}
  \fonte{os autores (2026).}
\end{table}

\section{Outra seção}
A acidose metabólica\index{Acidose!metabólica} é discutida aqui como exemplo de subentrada no índice.

\begin{referencias}
\item SOBRENOME, Nome; SOBRENOME, Nome. Título do artigo. \textbf{Nome da Revista}, v. 0, n. 0, p. 00-00, 2024. DOI: https://doi.org/10.0000/xxxx.
\item SOBRENOME, Nome. \textbf{Título do livro}. 2. ed. Cidade: Editora, 2023.
\end{referencias}

\capitulo[Nome da Autora]{Título do segundo capítulo}

\begin{resumo}
Resumo do segundo capítulo.
\end{resumo}
\palavraschave{termo; termo; termo.}

\section{Seção}
Cada capítulo abre em página ímpar. A hemogasometria\index{Hemogasometria} também aparece aqui, e o índice remissivo reúne as
páginas automaticamente. Em cães\index{Cães} e gatos\index{Gatos}, os valores de referência diferem.

\section{Considerações finais}
Síntese e aplicações práticas do conteúdo.

\begin{referencias}
\item SOBRENOME, Nome. \textbf{Título}. Cidade: Editora, 2022.
\end{referencias}


\backmatter
\chapter{Sobre os autores}
\sobreautor{foto-autor-exemplo}{Nome Completo do Autor}
  {Médico-veterinário, Doutor em Ciência Animal pela Universidade Exemplo. Professor da Instituição Exemplo, com atuação em clínica médica
   de pequenos animais e medicina laboratorial. Minicurrículo de até seis linhas, em terceira pessoa.}
  {ORCID 0000-0000-0000-0000 · Lattes http://lattes.cnpq.br/0000000000000000}
\sobreautor{foto-autor-exemplo}{Nome Completo da Autora}
  {Médica-veterinária, Mestra em Medicina Veterinária pela Universidade Exemplo, com atuação em patologia clínica.}
  {ORCID 0000-0000-0000-0000 · Lattes http://lattes.cnpq.br/0000000000000000}

\printindex
\colofao
\end{document}
